% !PRINTED 314
\weberpage{314}
Vergleichen wir nun hiermit den \gesperrt{Budan'schen} Satz, so haben wir die drei Functionen $f(x)$, $f'(x)$, $f''(x)$, also

% SOURCE FIGURE CROP: Bd. I registered scan 340 (printed p. 314), 600 dpi,
% crop x=300,y=700,w=2700,h=1440; source-faithful raster, Fig. 16--17.
\weberfigureimage{weber-b1-fig16-17.png}{0.94}

\weq{4}{x^2 + a\,x + b,\quad 2\,x + a,\quad 1}
in Bezug auf die Zeichenwechsel zu untersuchen.

Die beiden geraden Linien
\[
x^2 + a\,x + b = 0,\quad 2\,x + a = 0
\]
theilen die Ebene in drei Felder, in denen die Reihe (4) die in der Fig. 18 angedeuteten Zeichenwechsel bietet.

% SOURCE FIGURE CROP: Bd. I registered scan 340 (printed p. 314), 600 dpi,
% crop x=350,y=2820,w=2650,h=1450; source-faithful raster, Fig. 18--19.
\weberfigureimage{weber-b1-fig18-19.png}{0.94}

Nehmen wir wieder zwei Werthe $x = \alpha$ und $x = \beta$, so ergeben sich sechs Felder, in denen der Verlust an Zeichenwechseln in der Fig. 19 angegeben ist. In den mit 0 und 1 bezeichneten
%CONT>>
